\subsection{An unrestricted finite basis}\label{app:owner-finite-basis}

The unrestricted basis uses words of unary actions on a common-root
semilattice. Completeness follows from guard coverage and two-source
absorption.

\begin{theorem}[Finite basis for one owner]\label{owner:thm:finite-basis}
The equational theory of $\{\del,H_a\}$ has a basis of thirty-six
identities, each using at most six variables. The basis below is
complete over arbitrary algebras, finite or infinite. Consequently
the minimum possible variable bound of an equational basis lies
between five and six.
\end{theorem}

Here and throughout this subsection, write $h$ for $H_a$.
We give the basis explicitly. Fix a variable $r$ and use
the following term abbreviations:
\begin{gather*}
 p_r(u)=u\del r,\qquad
 u\vee_r v=(u\del r)\del(v\del r),\qquad 0=r,\\
 F_{ab}(u)=h(r,u\del a,b),\qquad
 G_a(u)=(u\del a)\del r,\qquad K_a(u)=h(r,u,a),\\
 c_a=a\del r,\qquad C_{ab}=h(r,a,b).
\end{gather*}
The subscript $r$ on the actions is suppressed. Juxtaposition of
actions means composition, with the rightmost action applied first.
In the action equations below, $x$ and $y$ abbreviate $u\del r$
and $v\del r$. Thus they range over $U_r=p_r(A)$ in any model.
Neither $0$ nor the actions are added symbols in the signature.
Every inequality $s\leq t$ abbreviates $s\vee_r t=t$.
Expanding these conventions gives ordinary unconditional identities.

The first six identities are (R1)--(R6). The next eleven are
\begin{align*}
 &\text{(S1)}\quad h(r,a,b)=h(r,b,a),\\
 &\text{(S2)}\quad h(r,a,a)=a\del r,\\
 &\text{(S3)}\quad p_r(h(r,a,b))=h(r,a,b),\\
 &\text{(S4)}\quad h(u\del r,a,b)=(u\del r)\del h(r,a,b),\\
 &\text{(S5)}\quad
 h(r,u\del a,v\del b)
   =h(r,u\del a,b)\vee_r h(r,a,v\del b),\\
 &\text{(S6)}\quad
 h(r,u\del a,b)=h(r,(u\del r)\del a,b),\\
 &\text{(S7)}\quad p_r(h(a,u,b))=G_a(h(r,u,b)),\\
 &\text{(S8)}\quad h(r,a,b)=h(r,r\del a,b),\\
 &\text{(S9)}\quad
 p_r((u\del a)\del(v\del a))
     =p_r(u\del a)\vee_r p_r(v\del a),\\
 &\text{(S10)}\quad u\del(v\del a)=(u\del a)\del(v\del a),\\
 &\text{(S11)}\quad p_r(u\del p_r(v))=u\del p_r(v).
\end{align*}
The remaining nineteen identities are the following action laws:
\begin{align*}
 &\text{(A1)}\quad p_r(F_{ab}(x))=F_{ab}(x),\\
 &\text{(A2)}\quad F_{ab}(x\vee_r y)=F_{ab}(x)\vee_r F_{ab}(y),\\
 &\text{(A3)}\quad F_{ab}(x)\leq F_{ab}(0)\vee_r x,\\
 &\text{(A4)}\quad F_{ab}F_{ab}(x)=F_{ab}(x),\\
 &\text{(A5)}\quad
 F_{ab}F_{cd}F_{ab}(x)=F_{ab}F_{cd}(x),\\
 &\text{(A6)}\quad F_{ab}G_a(x)=F_{ab}(x),\\
 &\text{(A7)}\quad F_{ab}K_a(x)=F_{ab}(x),\\
 &\text{(A8)}\quad F_{ab}K_b(x)=F_{ab}(x),\\
 &\text{(A9)}\quad F_{ab}(c_a)=F_{ab}(0),\\
 &\text{(A10)}\quad F_{ab}(c_b)=F_{ab}(0),\\
 &\text{(A11)}\quad G_aK_bF_{ab}(x)=G_aK_b(x),\\
 &\text{(A12)}\quad G_aF_{ab}(x)=G_aK_b(x),\\
 &\text{(A13)}\quad K_b(C_{ab})=K_b(c_a),\\
 &\text{(A14)}\quad C_{ab}\leq c_a\vee_r c_b,\\
 &\text{(A15)}\quad C_{ab}\leq K_b(c_a),\\
 &\text{(A16)}\quad G_a(x)=F_{aa}(x),\\
 &\text{(A17)}\quad K_a(x)=F_{ra}(x),\\
 &\text{(A18)}\quad G_r(x)=x,\\
 &\text{(A19)}\quad K_r(x)=x.
\end{align*}
All thirty-six identities hold in $\mathbf F$. This is a finite soundness
check, independent of the completeness proof below. The largest
variable set occurs in (A5): $r,a,b,c,d,u$, with $x=u\del r$.

\paragraph{Common-root words.}
Work now in an arbitrary model of these thirty-six identities.
Reset makes $U_r$ a join semilattice with bottom $r$, and its join
is the restriction of $\del$. For example, (R3) and (R6) give
\[
 (u\del r)\del(v\del r)=u\del(v\del r).
\]
Together with (R1) and (S11), this gives closure, commutativity,
and the bottom laws. For fixed points $x,y,z$ of $p_r$, (R6)
and closure give
\[
 (x\del y)\del z
 =z\del(x\del y)=x\del(z\del y)=x\del(y\del z),
\]
which proves associativity. We write $\vee$ for this join.
The actions preserve $U_r$ and finite nonempty joins, by (A1)--(A2),
(A16)--(A17). In particular they are monotone. An \emph{action word}
is a finite composition of the $F_{ab}$; the empty word is the identity.

Every term with root $r$ is derivably a finite nonempty join of
values $P(0)$ of action words. First, every term satisfies
$t=t\del\rho(t)$ derivably. The variable case is (R1).
For a reset term whose right child has root $j$, (S11) preserves
being fixed by $p_j$. For an $h$ term whose first child has root $i$,
(S4) writes it as a join of two elements of $U_i$, using (S3)
for the second element. This proves root preservation by induction.
More generally, put
$\bar t=p_r(t)$. If $s,t,u$ have roots $i,j,k$, respectively, the
normalization equations give
\begin{align}
 \overline{s\del t}&=G_j(\bar s)\vee\bar t,\label{owner:eq:word-reset}\\
 \overline{h(s,t,u)}&=\bar s\vee
 G_i\bigl(F_{jk}(\bar t)\vee F_{kj}(\bar u)\bigr).
 \label{owner:eq:word-h}
\end{align}
For \eqref{owner:eq:word-reset}, use $t=t\del j$, then
(S10), (S9), and (R5). For \eqref{owner:eq:word-h}, first use
(S4) to separate the first child, project with (S9), and use
(S7) to change its root to $r$. Equations (S5), (S1), and (S6)
then separate and project the other two children.
The initial case is $\bar a=G_a(0)$ by (R2).
Distribute the actions through joins and use (A16) to replace each
$G$ by an $F$. This proves the asserted normal form by induction.

\paragraph{Word identities.}
The following consequences are equational derivations, valid in
every model of the displayed basis. By (A4)--(A5), a repeated
primitive action can be removed even across an arbitrary word:
\[
 FQF=FQ.
\]
Indeed, for $Q=GV$, insert an $F$ after $G$, apply the induction
hypothesis to $FVF$, then remove the inserted $F$ using (A5).
The empty case is (A4). Removing the letters of a repeated word
one at a time therefore gives
\begin{equation}\label{owner:eq:word-band}
 PQP=PQ
 \qquad\text{for all action words }P,Q.
\end{equation}
Here and below equality of words means equality of their actions
on $U_r$. The $G$ and $K$ actions are also words by (A16)--(A17).

Every word is contractive up to its value at zero:
\begin{equation}\label{owner:eq:word-contract}
 P(x)\leq P(0)\vee x.
\end{equation}
For $P=FQ$, induction and (A2)--(A3) give
$FQ(x)\leq FQ(0)\vee F(x)\leq FQ(0)\vee F(0)\vee x$;
monotonicity absorbs $F(0)$ into $FQ(0)$.

For a word $P$, let $M(P)$ contain $r$ and the first indices of
its letters, and let $N(P)$ contain $r$ and both indices of its
letters. In particular $M(P)\subseteq N(P)$. Equations
(A6)--(A10), (A18)--(A19), and \eqref{owner:eq:word-band} imply
\begin{align}
 PG_a=P&\quad(a\in M(P)),&
 PK_b=P&\quad(b\in N(P)),\label{owner:eq:word-tests}\\
 P(c_b)&=P(0)&&\quad(b\in N(P)).\label{owner:eq:word-constants}
\end{align}
For example, if $P$ contains $F_{ab}$, then $PF_{ab}=P$.
Appending $G_a$, $K_a$, or $K_b$ and using (A6)--(A8) proves
the corresponding equations. Applying the same equality to $c_a$
or $c_b$ and using (A9)--(A10) proves the constant equations.
The index $r$ follows from the bottom laws and (A18)--(A19).

Two further consequences will remove a witness word:
\begin{align}
 PF_{ab}&=P&& (a\in M(P),\ b\in N(P)),
                 \label{owner:eq:word-capture}\\
 PF_{ab}&=PK_b&& (a\in M(P)).
                 \label{owner:eq:word-weak}
\end{align}
For the first, insert $G_aK_b$ after $P$ using
\eqref{owner:eq:word-tests}, then apply (A11).
For the second insert $G_a$ and apply (A12).
Finally, (A13), (S2), and \eqref{owner:eq:word-tests} give
\begin{equation}\label{owner:eq:word-terminal}
 P(C_{jt})=P(c_j)\qquad(t\in N(P)\cup\{j\}).
\end{equation}

\paragraph{Guards and first occurrences.}
The semantic input is the already proved guard completeness theorem.
For a word value $PF_{ab}(0)=P(C_{ab})$, the new source-$a$ guard
has ambiguity tests $M(P)$ and passing tests $N(P)\cup\{b\}$;
the source-$b$ guard has passing tests $N(P)\cup\{a\}$.
Duplicate passing tests on ambiguous inputs are harmless.
All word values also have the unconditional root guard.
This description follows directly from the guard recursion:
$F_{ab}$ has multiplier $M_aN_b$ and constant contributions
$V_aN_b$ and $V_bN_a$, where $V_j$ denotes the source marker.

Only the first occurrence of a source in a word is needed as a
covering witness. If $j$ occurs first as the first index of an
earlier letter, its companion's passing test appears among the
later passing tests. If it occurs as the second index, its
companion's passing test is implied by a later ambiguity test.
In both cases the earlier source-$j$ guard covers every later
source-$j$ guard in that word. Thus a covering witness can always
be chosen in the form $Q(C_{jt})$, with neither index of any
letter of $Q$ equal to $j$. Such a witness is below its full word
value, by monotonicity and $0\leq R(0)$ for the omitted suffix $R$.

\begin{lemma}[Two-source absorption]\label{owner:lem:two-source}
Let $W$ be a join of word values. Suppose $P(0)\leq W$ is already
derivable, and each of the two new source guards of $P(C_{ab})$
is covered by a guard of $W$. Then $P(C_{ab})\leq W$ is derivable
from the thirty-six identities.
\end{lemma}
\begin{proof}
For a source $j\in\{a,b\}$ let $k$ be its companion. If
$j\in N(P)$, then $P(c_j)=P(0)\leq W$ by
\eqref{owner:eq:word-constants}. Otherwise choose a first-occurrence
witness $Q(C_{jt})\leq W$. Guard coverage says that every letter
$F_{uv}$ of $Q$ satisfies
\[
 u\in M(P),\qquad v\in N(P)\cup\{k\},
 \qquad t\in N(P)\cup\{k,j\}.
\]
The source $j$ does not occur in $Q$, which removes the otherwise
allowed passing test on $j$ itself.

Equations \eqref{owner:eq:word-capture}--\eqref{owner:eq:word-weak}
reduce $PQ$ to either $P$ or $PK_k$. The latter case is needed
only if $k\notin N(P)$ and $Q$ contains a letter with second
index $k$. This reduction proceeds one letter at a time. After
$PK_k$ has appeared, insert $P$ using $PK_kP=PK_k$ before
reducing the next letter; (A4) and
\eqref{owner:eq:word-band} remove repeated actions.

The witness can be prefixed without leaving $W$:
\[
 P Q(C_{jt})\leq P(0)\vee Q(C_{jt})\leq W,
\]
by \eqref{owner:eq:word-contract}.
If $PQ=P$ and $t\in N(P)\cup\{j\}$, this proves
$P(c_j)\leq W$ using \eqref{owner:eq:word-terminal}.
If $PQ=P$ and the remaining possibility $t=k$ occurs, it
already proves $P(C_{jk})\leq W$.

If $PQ=PK_k$, then
\[
 PK_k(C_{jt})=PK_k(c_j).
\]
For $t=k$ this is (A13). For
$t\in N(P)\cup\{j\}$, insert $P$ after $K_k$, apply
\eqref{owner:eq:word-terminal}, and remove $P$ again.
Equation (A15) now gives
$P(C_{jk})\leq PK_k(c_j)\leq W$.

Thus each source either supplies the entire pair value or proves
$P(c_j)\leq W$. In the former case the lemma is proved, using
(S1) if the source was $b$. If both sources supply only their
own constants, (A14) and preservation of joins give
\[
 P(C_{ab})\leq P(c_a)\vee P(c_b)\leq W.
\]
This also covers equal indices and indices already in $N(P)$.
\end{proof}

\begin{proof}[Proof of \cref{owner:thm:finite-basis}]
Soundness was checked on $\mathbf F$, which generates the
one-owner theory. For completeness, normalize a valid identity
$s=t$. Its two terms have the same root $r$, as is seen on the
projection subalgebra $\{A,B\}$. We obtain joins $W,V$ of word
values on $U_r$.

Take any word $F_1\cdots F_m$ occurring in $V$. Starting with
$0\leq W$, prove successively that each prefix value
$F_1\cdots F_i(0)$ is at most $W$. At the next letter, the two
source guards are guards of the full word, hence guards of $V$.
Since $V=W$ is valid in $\mathbf F$, guard completeness supplies
covering guards in $W$. Apply \cref{owner:lem:two-source}.
The resulting inequality for the whole word, and then for their
join, proves $V\leq W$. Reversing the roles proves $W\leq V$.
All inequalities used here abbreviate derivable join identities,
so semilattice antisymmetry completes an equational derivation
of $s=t$. No conservativity or cardinality assumption was made.
The lower variable bound is \cref{owner:thm:no-four-variable-basis}.
\end{proof}
